% Noether P06 Korean producer draft — UNCHECKED
% T10-U142; ED0001 lines5380--5387; source 903 LF B / 8A0DBEF3D8275FD6A50BE9876AAC610FFC6D1CA8E186DABC7504AB947687486F
% Pointer v008 / ED0001: F13E6D896DE6403829FE902609668AB3E3FCA8C3C7FAA07BE5F7A7A72A4C33D8 / D1F06B311F6CBD991DD247D745DD9A72DDE326A20396DF43CFE0C8EDB1593CDB
% Translation only; no review/build/render/approval. Same-smallest-field proof, common-divisor possibility, footnote example, and every coefficient/exponent remain checker holds.

$\Jdom_{n\rho}$와 $\Gdom_{n\rho}$는 이를 포함하는 동일한 최소 체 $\Kfield_{n\rho}$를 갖기 때문이다. 또한 $\Jdom_{n\rho}$의 모든 인볼루션형식은 $\Kfield_{n\rho}$의 한 인볼루션형식에 $\Jdom_{n\rho}$의 다항식을 곱하여 생긴다. 따라서 그 계수들은 $\Gdom_{n\rho}$에 속하지만, $\Gdom_{n\rho}$에서 공약수를 가질 수 있다.\srcfn{*)}{예를 들어 $x_1^2$과 $x_1x_2$의 모든 정유리 결합으로 이루어진 정역 $\Jdom_{22}$에 대해서는 다음 인볼루션형식이 얻어진다:
\[
\Psi(z,u)=x_1^2z^2-x_1^2u_1^2-2x_1^2x_2u_1u_2-x_1^2x_2^2u_2^2.
\]
$x_1^2$은 $\Jdom_{22}$에, 따라서 더더욱 이를 포함하는 최소 상대정역 $\Gdom_{22}$에 속하고, 그러므로 $\Gdom_{22}$에서 공약수이다. 이에 따라 $\Gdom_{22}$의 인볼루션형식으로 다음이 얻어진다:
\[
\Psi(z,u)=z^2-x_1^2u_1^2-2x_1x_2u_1u_2-x_2^2u_2^2.
\]}
